% FIGURE NEEDED: sketch of Pluto (P) and Charon (C) on the line joining
% their centers, with the center of mass CM marked between them, unit vector
% r-hat pointing from Charon toward Pluto, points A and B on Charon's
% near/far side, and the rotation arrows; 2021 T9 solution PDF, page 1.
Let P and C be the centers of Pluto and Charon respectively, and CM their
center of mass. Clearly $\overline{CM - P} = \frac{1}{9}R$ and
$\overline{CM - C} = \frac{8}{9}R$. The effective gravitational field at
points A and B are given by:
\begin{center}
[By recognizing the vectorial nature of gravities/forces/fields]
\hfill [4 points]
\end{center}
\begin{gather*}
g_A = \frac{GM_P}{(R-r)^2}\hat{r} - \frac{GM_C}{r^2}\hat{r}
  - \omega^2\left(\tfrac{8}{9}R - r\right)\hat{r} \tag*{[1 point]} \\
g_B = \frac{GM_P}{(R+r)^2}\hat{r} + \frac{GM_C}{r^2}\hat{r}
  - \omega^2\left(\tfrac{8}{9}R + r\right)\hat{r} \tag*{[1 point]}
\end{gather*}
where $\omega^2 = 9\dfrac{GM_C}{R^3}$ is the angular speed common to both
Pluto and Charon. Thus, the gravitational accelerations at these points are:
\begin{center}
[by finding angular speed] [2 points]
\end{center}
\begin{gather*}
g_A = \frac{GM_P}{(R-r)^2} - \frac{GM_C}{r^2}
  - \frac{GM_C}{R^2}\left(8 - 9\frac{r}{R}\right) \\
g_B = \frac{GM_P}{(R+r)^2} + \frac{GM_C}{r^2}
  - \frac{GM_C}{R^2}\left(8 + 9\frac{r}{R}\right)
\end{gather*}
Factorizing the expression to obtain formulas proportional to Charon's mass
[1 point]:
\begin{gather*}
g_A = -\frac{GM_C}{r^2}\left[1
  - \frac{M_P}{M_C}\left(\frac{r}{R}\right)^2
    \frac{1}{\left(1 - \frac{r}{R}\right)^2}
  + \left(\frac{r}{R}\right)^2\left(8 - 9\frac{r}{R}\right)\right] \\
g_B = \frac{GM_C}{r^2}\left[1
  + \frac{M_P}{M_C}\left(\frac{r}{R}\right)^2
    \frac{1}{\left(1 + \frac{r}{R}\right)^2}
  - \left(\frac{r}{R}\right)^2\left(8 + 9\frac{r}{R}\right)\right]
\end{gather*}
Replacing the numerical value of the ratio of the two masses:
\begin{gather*}
g_A = -\frac{GM_C}{r^2}\left[1
  - \left(\frac{r}{R}\right)^2\frac{8}{\left(1 - \frac{r}{R}\right)^2}
  + \left(\frac{r}{R}\right)^2\left(8 - 9\frac{r}{R}\right)\right] \\
g_B = \frac{GM_C}{r^2}\left[1
  + \left(\frac{r}{R}\right)^2\frac{8}{\left(1 + \frac{r}{R}\right)^2}
  - \left(\frac{r}{R}\right)^2\left(8 + 9\frac{r}{R}\right)\right]
\end{gather*}
Noting that the expressions in square brackets are dimensionless, it is
helpful to replace the given values of $r$ and $R$, resulting [2 points by
expressions of $g_A$ and $g_B$]:
\begin{gather*}
g_A = -\frac{GM_C}{r^2}\,(0{,}99929) \\
g_B = \frac{GM_C}{r^2}\,(0{,}99933)
\end{gather*}
The percentage difference with respect to Charon's normal gravity
$g_0 = \dfrac{GM_C}{r^2}$ is
\begin{equation*}
\frac{\big||g_A| - |g_B|\big|}{g_0}\cdot 100\% = 0{,}004\%
  = 4 \times 10^{-3}\,\% \tag*{[1 point]}
\end{equation*}
% note: the source (Bogota) uses decimal commas ("0,99929", "0,004%");
% kept as printed.
